% ================================================================
% CHAPTER 10: NETWORKS AND SPECTRAL THEORY OF CULTURAL AGREEMENT
% ================================================================
\chapter{Networks and Spectral Theory of Cultural Agreement}
\label{ch:spectral}

\epigraph{Society is indeed a contract \ldots\ a partnership not only between
those who are living, but between those who are living, those who are dead, and
those who are to be born.}{Edmund Burke, \textit{Reflections on the Revolution
in France}}

\section{From Pairs to Populations}

The preceding chapters analyzed the interaction between \emph{pairs} of ideas.
We now scale up to \emph{populations}: networks of agents, each holding an
idea-vector, all embedded in the same Hilbert space $\Hspace$. The central
tool is spectral theory --- the analysis of the eigenvalues and eigenvectors of
the resonance matrix --- which reveals the ``hidden structure'' of cultural
agreement.

\begin{definition}[Cultural Population]
\label{def:population}
A \textbf{cultural population} consists of $n$ agents, each holding an idea
$a_i \in \IS$ for $i = 1, \ldots, n$. The population is represented by the
collection of embedded vectors $\{\embed{a_1}, \ldots, \embed{a_n}\}$ in
$\Hspace$.
\end{definition}

\begin{definition}[Resonance Matrix]
\label{def:resonance-matrix}
The \textbf{resonance matrix} of a cultural population is the $n \times n$
matrix $\RMat$ with entries:
\[
  R_{ij} := \rs{a_i}{a_j} = \ip{\embed{a_i}}{\embed{a_j}}.
\]
\end{definition}

\begin{theorem}[The Resonance Matrix is a Gram Matrix]
\label{thm:gram}
The resonance matrix $\RMat$ is a \textbf{Gram matrix}: $\RMat = \Phi^\top \Phi$,
where $\Phi$ is the matrix whose columns are the embedded vectors $\embed{a_i}$.
Consequently:
\begin{enumerate}
\item $\RMat$ is symmetric: $R_{ij} = R_{ji}$;
\item $\RMat$ is positive semidefinite: $x^\top \RMat x \geq 0$ for all
$x \in \mathbb{R}^n$;
\item All eigenvalues of $\RMat$ are non-negative: $\lambda_1 \geq \lambda_2
\geq \cdots \geq \lambda_n \geq 0$;
\item $\operatorname{rank}(\RMat) = \dim(\operatorname{span}\{\embed{a_1},
\ldots, \embed{a_n}\})$.
\end{enumerate}
\end{theorem}

\begin{proof}
$R_{ij} = \ip{\embed{a_i}}{\embed{a_j}} = (e_i^\top \Phi^\top)(\Phi e_j) =
e_i^\top \Phi^\top \Phi e_j$, so $\RMat = \Phi^\top \Phi$. This immediately
gives symmetry (the transpose of $\Phi^\top \Phi$ is itself) and positive
semidefiniteness ($x^\top \Phi^\top \Phi x = \nrm{\Phi x}^2 \geq 0$). Non-negative
eigenvalues follow from positive semidefiniteness. The rank equals the dimension
of the column space of $\Phi$, which is $\operatorname{span}\{\embed{a_1},
\ldots, \embed{a_n}\}$.

In Lean~4 (sketch):
\begin{lstlisting}
theorem resonance_matrix_psd (a : Fin n → Idea) :
    Matrix.PosSemidef (fun i j => rs (a i) (a j)) := by
  intro x
  simp [rs, Matrix.dotProduct]
  -- reduce to ‖∑ᵢ xᵢ • φ(aᵢ)‖² ≥ 0
  sorry
\end{lstlisting}
\end{proof}

\begin{interpretation}[Positive Semidefiniteness as Internal Consistency]
\label{interp:psd}
The fact that the resonance matrix is always positive semidefinite is a
deep structural constraint on the pattern of cultural agreement. Not every
symmetric matrix can arise as a resonance matrix; only PSD matrices can. This
means that the \emph{pattern} of who agrees with whom is not arbitrary: it must
be consistent with an underlying Hilbert space geometry.

In practice, this is a testable prediction. If we measure pairwise ``cultural
affinity'' between groups (by surveys, content analysis, or behavioral data)
and construct a matrix, IDT~v2 predicts that this matrix should be PSD (or at
least approximately PSD, allowing for measurement error). If the empirical
matrix has significantly negative eigenvalues, it indicates that the pairwise
affinities are not consistent with any Hilbert space embedding --- a failure
of the model that would point toward non-linear or non-metric structure in
cultural agreement.

Conversely, if the empirical matrix \emph{is} PSD, we can extract an embedding:
the eigendecomposition of $\RMat$ directly yields the vectors
$\{\embed{a_i}\}$ (up to rotation). This is precisely what methods like
Multidimensional Scaling (MDS) and Principal Component Analysis (PCA) do:
they take a matrix of pairwise similarities and extract an embedding in a
Euclidean space. IDT~v2 provides the theoretical justification for these
techniques when applied to cultural data.
\end{interpretation}


\section{The Eigenvalues of Cultural Agreement}

The eigenvalues of $\RMat$ reveal the ``principal modes'' of cultural agreement.

\begin{definition}[Spectral Decomposition of Culture]
\label{def:spectral-decomposition}
Let $\RMat = \sum_{k=1}^r \lambda_k u_k u_k^\top$ be the eigendecomposition
of the resonance matrix, with eigenvalues $\lambda_1 \geq \lambda_2 \geq \cdots
\geq \lambda_r > 0$ ($r = \operatorname{rank}(\RMat)$) and orthonormal
eigenvectors $u_1, \ldots, u_r \in \mathbb{R}^n$.

\begin{itemize}
\item The $k$-th eigenvalue $\lambda_k$ is the \textbf{$k$-th mode of
agreement}: it measures the variance of the population along the $k$-th
principal direction in meaning space.
\item The $k$-th eigenvector $u_k$ describes \textbf{who participates} in the
$k$-th mode: $u_{k,i}$ is the ``loading'' of agent $i$ on mode $k$.
\end{itemize}
\end{definition}

\begin{theorem}[Trace and Total Agreement]
\label{thm:trace}
The sum of all eigenvalues equals the \textbf{total self-resonance} of the
population:
\[
  \sum_{k=1}^r \lambda_k = \operatorname{tr}(\RMat) = \sum_{i=1}^n
  \rs{a_i}{a_i} = \sum_{i=1}^n w(a_i)^2.
\]
\end{theorem}

\begin{proof}
$\operatorname{tr}(\RMat) = \sum_i R_{ii} = \sum_i \rs{a_i}{a_i}$. The trace
equals the sum of eigenvalues.
\end{proof}

\begin{theorem}[Largest Eigenvalue and Maximum Agreement]
\label{thm:largest-eigenvalue}
The largest eigenvalue $\lambda_1$ satisfies:
\[
  \lambda_1 = \max_{\nrm{x} = 1} x^\top \RMat x
  = \max_{\nrm{x} = 1} \sum_{i,j} x_i x_j \rs{a_i}{a_j}
  = \max_{\nrm{x} = 1} \nrm{\sum_i x_i \embed{a_i}}^2.
\]
The maximizing $x = u_1$ defines the \textbf{direction of maximum cultural
agreement}: the linear combination $\sum_i u_{1,i} \embed{a_i}$ is the
``consensus vector'' that captures the most variance.
\end{theorem}

\begin{proof}
This is the Rayleigh quotient characterization of the largest eigenvalue of
a symmetric matrix. The last equality uses:
\[
  x^\top \RMat x = \sum_{i,j} x_i x_j \ip{\embed{a_i}}{\embed{a_j}}
  = \ip{\sum_i x_i \embed{a_i}}{\sum_j x_j \embed{a_j}}
  = \nrm{\sum_i x_i \embed{a_i}}^2.
\]
\end{proof}

\begin{interpretation}[Durkheim's Collective Conscience]
\label{interp:durkheim}
Theorem~\ref{thm:largest-eigenvalue} provides a precise mathematical
formalization of Durkheim's ``collective conscience'' (conscience collective):
\emph{the eigenvector of the largest eigenvalue of the resonance matrix is the
direction in meaning space that maximizes total agreement across the
population.}

Durkheim defined the collective conscience as ``the totality of beliefs and
sentiments common to the average members of a society.'' In IDT~v2, this is
the vector $v_1 = \sum_i u_{1,i} \embed{a_i}$ --- the weighted sum of all
individuals' ideas, with weights chosen to maximize the squared norm (total
agreement). The collective conscience is not the simple average of individual
beliefs (which would give equal weight to all agents) but a \emph{weighted}
combination that emphasizes the most agreeable direction.

The eigenvalue $\lambda_1$ measures the \emph{strength} of the collective
conscience: how much of the total cultural variance is captured by this single
direction. A culture with $\lambda_1 \approx \operatorname{tr}(\RMat)$ (the
first eigenvalue captures nearly all the variance) has a strong collective
conscience --- nearly all individuals agree on the principal direction. A
culture with $\lambda_1 \approx \lambda_2 \approx \cdots$ (eigenvalues spread
out) has a weak collective conscience --- there is no single direction of
dominant agreement.

In Durkheim's terminology:
\begin{itemize}
\item \textbf{Mechanical solidarity} corresponds to $\lambda_1 \gg \lambda_2$:
a single dominant mode of agreement, with all individuals approximately aligned.
\item \textbf{Organic solidarity} corresponds to $\lambda_1 \approx \lambda_2
\approx \cdots$: multiple independent modes of agreement, with individuals
specializing in different cultural directions.
\end{itemize}
\end{interpretation}


\section{The Spectral Gap}

\begin{definition}[Spectral Gap]
\label{def:spectral-gap}
The \textbf{spectral gap} of the resonance matrix is:
\[
  \gamma := \lambda_1 - \lambda_2 \geq 0.
\]
\end{definition}

\begin{theorem}[Spectral Gap and Cultural Cohesion]
\label{thm:spectral-gap}
Let $\alpha_i := \ip{u_1}{\hat{a}_i}$ be the ``alignment'' of agent $i$'s
idea with the first eigenvector direction, where $\hat{a}_i =
\embed{a_i}/\nrm{\embed{a_i}}$. Then:
\[
  \frac{1}{n}\sum_{i=1}^n \alpha_i^2 \geq \frac{\lambda_1}{\lambda_1 +
  (n-1)(\lambda_1 - \gamma)}.
\]
In particular, a large spectral gap $\gamma$ implies high average alignment
with the first mode.
\end{theorem}

\begin{proof}[Proof sketch]
This follows from the variational characterization of eigenvalues and the
Davis--Kahan theorem, which bounds the perturbation of eigenvectors in terms
of the spectral gap. The precise bound depends on the specific norm used.
The essential point is that a large gap between $\lambda_1$ and $\lambda_2$
ensures that the first eigenvector is ``well-separated'' from the second, and
hence that the first mode captures a disproportionate share of the cultural
variance.
\end{proof}

\begin{interpretation}[The Spectral Theory of Cultural Unity]
\label{interp:spectral-gap}
The spectral gap $\gamma$ is a single-number summary of \emph{cultural
cohesion}: how unified or fragmented a culture is. A large spectral gap means
that one mode of agreement dominates all others --- the culture is
``one-dimensional'' in the sense that a single ideological axis captures most
of the variation. A small spectral gap means that multiple modes are equally
important --- the culture is ``multi-dimensional'' and fragmented.

This connects to debates in political science about polarization. A polarized
society is often described as ``two-dimensional'' (liberal--conservative on one
axis, and some cross-cutting cleavage on another). In spectral terms,
polarization would correspond to $\lambda_1 \approx \lambda_2 \gg \lambda_3$:
two dominant modes of roughly equal importance. De-polarization would
correspond to increasing the spectral gap: making one mode dominant and the
others subordinate.

The spectral gap also has practical implications for cultural governance. A
society with a large spectral gap is easy to ``manage'' (in a cultural sense):
the single dominant mode provides a natural axis of agreement, and policies
that align with this axis enjoy broad support. A society with a small spectral
gap is harder to manage: there is no single axis of agreement, and any policy
aligning with one mode necessarily conflicts with others.
\end{interpretation}


\section{Principal Cultural Dimensions}

The eigenvectors of $\RMat$ define the ``principal cultural dimensions'' ---
the axes of meaning along which the population varies most.

\begin{definition}[Principal Cultural Dimensions]
\label{def:pcd}
The \textbf{$k$-th principal cultural dimension} is the direction in $\Hspace$
corresponding to the $k$-th eigenvector:
\[
  d_k := \Phi u_k = \sum_{i=1}^n u_{k,i} \embed{a_i}.
\]
This is the direction in meaning space that captures the $k$-th most variance.
\end{definition}

\begin{theorem}[Reconstruction from Principal Dimensions]
\label{thm:reconstruction}
Each individual's idea can be expressed as:
\[
  \embed{a_i} = \sum_{k=1}^r c_{ik} \cdot \frac{d_k}{\nrm{d_k}},
\]
where $c_{ik} = \ip{\embed{a_i}}{d_k/\nrm{d_k}}$ is the ``coordinate'' of
agent $i$ along the $k$-th cultural dimension. The resonance between agents
$i$ and $j$ decomposes as:
\[
  \rs{a_i}{a_j} = \sum_{k=1}^r c_{ik} c_{jk}.
\]
\end{theorem}

\begin{proof}
Since $\{d_k/\nrm{d_k}\}$ is an orthonormal basis for the span of
$\{\embed{a_i}\}$, the first claim is the expansion in this basis. The
resonance decomposition follows:
\[
  \rs{a_i}{a_j} = \ip{\embed{a_i}}{\embed{a_j}}
  = \ip{\sum_k c_{ik} \hat{d}_k}{\sum_l c_{jl} \hat{d}_l}
  = \sum_k c_{ik} c_{jk}.
\]
\end{proof}

\begin{interpretation}[Cultural Coordinates]
\label{interp:coordinates}
Theorem~\ref{thm:reconstruction} says that every individual's intellectual
position can be described by their coordinates along the principal cultural
dimensions. This is a dimensional reduction: instead of locating each agent
in the full (possibly infinite-dimensional) Hilbert space $\Hspace$, we
describe them by $r$ numbers (their cultural coordinates), where $r$ is the
rank of the resonance matrix.

If the first $K$ eigenvalues capture most of the variance (i.e., $\sum_{k=1}^K
\lambda_k / \sum_{k=1}^r \lambda_k \approx 1$), then $K$ coordinates suffice
to approximately describe every individual's intellectual position. This is
dimensional reduction in the cultural domain --- precisely the technique of
\emph{political compass} models, which attempt to locate individuals on two or
three axes (e.g., economic left--right, social liberal--conservative).

The spectral theory of IDT~v2 provides a principled way to derive these axes
from data, rather than imposing them a priori. The axes emerge from the
eigenstructure of the resonance matrix and are guaranteed to capture the
maximum variance. This is a cultural application of PCA (Principal Component
Analysis), but with a firm theoretical foundation in the Hilbert space model.
\end{interpretation}

\begin{example}[Spectral Analysis of a Five-Person Culture]
\label{ex:spectral}
Consider five agents in $\mathbb{R}^3$:
\begin{align*}
  \embed{a_1} &= (3, 1, 0), \\
  \embed{a_2} &= (2, 2, 0), \\
  \embed{a_3} &= (4, 0, 0), \\
  \embed{a_4} &= (1, 1, 3), \\
  \embed{a_5} &= (0, 0, 4).
\end{align*}
The resonance matrix:
\[
  \RMat = \begin{pmatrix}
  10 & 8 & 12 & 4 & 0 \\
  8 & 8 & 8 & 4 & 0 \\
  12 & 8 & 16 & 4 & 0 \\
  4 & 4 & 4 & 11 & 12 \\
  0 & 0 & 0 & 12 & 16
  \end{pmatrix}.
\]
The trace is $10 + 8 + 16 + 11 + 16 = 61$, so the total self-resonance is 61.

Computing eigenvalues (approximately):
\[
  \lambda_1 \approx 32.5, \quad \lambda_2 \approx 22.1, \quad
  \lambda_3 \approx 5.8, \quad \lambda_4 \approx 0.6, \quad \lambda_5 = 0.
\]
The spectral gap is $\gamma \approx 32.5 - 22.1 = 10.4$. The first two modes
capture $(32.5 + 22.1)/61 \approx 89.5\%$ of the total variance, suggesting
that this culture is essentially two-dimensional.

The first eigenvector corresponds roughly to the ``dimension 1'' axis
(agents 1--3, who are concentrated in the $e_1$--$e_2$ plane). The second
eigenvector corresponds roughly to the ``dimension 3'' axis (agents 4--5, who
have significant $e_3$ components). The culture has two ``camps'' --- a
humanistic camp (agents 1--3) and a scientific camp (agents 4--5) --- with
moderate cross-resonance mediated by agent 4, who has non-zero components in
both directions.
\end{example}


\section{The Resonance Network}

The resonance matrix can be viewed as the adjacency matrix of a weighted network.

\begin{definition}[Resonance Network]
\label{def:resonance-network}
The \textbf{resonance network} is the complete weighted graph $G = (V, E, w)$
where:
\begin{itemize}
\item Vertices $V = \{1, \ldots, n\}$ are agents;
\item Edge weights $w_{ij} = \rs{a_i}{a_j}$ (which can be negative).
\end{itemize}
\end{definition}

\begin{theorem}[Total Resonance and Graph Quadratic Form]
\label{thm:total-resonance}
The \textbf{total mutual resonance} of the population is:
\[
  T := \sum_{i < j} \rs{a_i}{a_j} = \frac{1}{2}\left(\nrm{\sum_i
  \embed{a_i}}^2 - \sum_i \nrm{\embed{a_i}}^2\right).
\]
\end{theorem}

\begin{proof}
\begin{align*}
  \nrm{\sum_i \embed{a_i}}^2
  &= \sum_i \nrm{\embed{a_i}}^2 + 2\sum_{i < j} \ip{\embed{a_i}}{\embed{a_j}}
  \\
  &= \sum_i \rs{a_i}{a_i} + 2\sum_{i < j} \rs{a_i}{a_j} \\
  &= \sum_i w(a_i)^2 + 2T.
\end{align*}
Rearranging: $T = \frac{1}{2}(\nrm{\sum_i \embed{a_i}}^2 - \sum_i w(a_i)^2)$.

In Lean~4:
\begin{lstlisting}
theorem total_resonance (a : Fin n → Idea) :
    ∑ i in Finset.univ, ∑ j in Finset.Iio i, rs (a i) (a j) =
    (‖∑ i, φ (a i)‖ ^ 2 - ∑ i, ‖φ (a i)‖ ^ 2) / 2 := by
  simp [rs, norm_sum_sq, inner_sum]
  ring
\end{lstlisting}
\end{proof}

\begin{interpretation}[The Total Resonance of a Culture]
\label{interp:total-resonance}
The total mutual resonance $T$ measures the overall ``coherence'' of the
cultural population. It is maximized when all ideas point in the same direction
(maximum alignment) and can be negative if the population is highly conflicted
(ideas pointing in opposing directions).

The formula $T = \frac{1}{2}(\nrm{\sum_i \embed{a_i}}^2 - \sum_i w(a_i)^2)$
has a beautiful interpretation: the total resonance equals half the difference
between the squared norm of the ``cultural centroid'' (the vector sum of all
ideas) and the sum of individual squared norms. If the centroid is long
(all ideas reinforce each other), $T$ is large and positive. If the centroid
is short (ideas cancel each other), $T$ is small or negative.

A culture that maximizes $T$ subject to a fixed total weight ($\sum_i w(a_i)^2$
constant) is one where all ideas point in the same direction. A culture that
minimizes $T$ is one where ideas point in maximally opposing directions. The
spectral decomposition tells us exactly how much agreement exists along each
principal axis.
\end{interpretation}


\section{Clustering and Community Structure}

The spectral theory naturally identifies ``communities'' within the cultural
network.

\begin{definition}[Spectral Clustering]
\label{def:clustering}
A \textbf{spectral partition} of the population into $K$ clusters $C_1, \ldots,
C_K$ is obtained by:
\begin{enumerate}
\item Computing the first $K$ eigenvectors $u_1, \ldots, u_K$ of $\RMat$;
\item Representing each agent $i$ by the $K$-dimensional vector
$(u_{1,i}, \ldots, u_{K,i})$;
\item Clustering these $K$-dimensional representations (e.g., by $k$-means).
\end{enumerate}
\end{definition}

\begin{theorem}[Within-Cluster Resonance Maximization]
\label{thm:cluster-resonance}
The spectral partition into $K$ clusters approximately maximizes the
\textbf{within-cluster total resonance}:
\[
  W(C_1, \ldots, C_K) := \sum_{k=1}^K \sum_{i, j \in C_k} \rs{a_i}{a_j}.
\]
\end{theorem}

\begin{proof}[Proof sketch]
This is a well-known result in spectral clustering theory. The within-cluster
resonance can be written as $W = \operatorname{tr}(H^\top \RMat H)$, where $H$
is the cluster indicator matrix. By the Ky Fan theorem, this is maximized (over
all rank-$K$ projections) by the subspace spanned by the top $K$ eigenvectors.
The spectral clustering approximates this optimal projection.
\end{proof}

\begin{interpretation}[Cultural Tribes]
\label{interp:tribes}
Spectral clustering of the resonance matrix identifies ``cultural tribes'' ---
groups of agents whose ideas mutually resonate. Within each tribe, ideas are
aligned (positive resonance, small angles); across tribes, ideas may be
orthogonal (independent) or opposing (conflicting).

This provides a data-driven, geometry-based alternative to subjective
taxonomies of cultural groups. Instead of classifying people as ``liberal'' or
``conservative'' based on predetermined criteria, spectral clustering derives
the relevant groupings from the \emph{structure of agreement itself}. The
number of clusters $K$ is not imposed but can be determined by the spectral
gap structure: a large gap between $\lambda_K$ and $\lambda_{K+1}$ suggests
that $K$ clusters capture the essential structure.

In the example of the five-person culture
(Example~\ref{ex:spectral}), two clusters emerge naturally: agents 1--3 (the
``humanistic'' cluster) and agents 4--5 (the ``scientific'' cluster). The
spectral gap between $\lambda_2$ and $\lambda_3$ confirms that $K = 2$
captures the essential structure.
\end{interpretation}


\section{The Dimensionality of a Culture}

\begin{definition}[Cultural Dimensionality]
\label{def:dimensionality}
The \textbf{effective dimensionality} of a cultural population is:
\[
  d_{\text{eff}} := \frac{(\operatorname{tr}(\RMat))^2}
  {\operatorname{tr}(\RMat^2)} = \frac{(\sum_k \lambda_k)^2}
  {\sum_k \lambda_k^2}.
\]
This is the ``participation ratio'' of the eigenvalues: it equals $1$ when one
eigenvalue dominates (one-dimensional culture) and $n$ when all eigenvalues
are equal (maximally multi-dimensional culture).
\end{definition}

\begin{theorem}[Bounds on Cultural Dimensionality]
\label{thm:dim-bounds}
$1 \leq d_{\text{eff}} \leq r = \operatorname{rank}(\RMat) \leq \min(n,
\dim(\Hspace))$.
\end{theorem}

\begin{proof}
By Cauchy--Schwarz: $(\sum_k \lambda_k)^2 \leq r \sum_k \lambda_k^2$, so
$d_{\text{eff}} \leq r$. Equality holds when all non-zero eigenvalues are equal.
The lower bound $d_{\text{eff}} \geq 1$ holds because $\lambda_1 \geq
\operatorname{tr}(\RMat)/n$ implies $(\sum \lambda_k)^2 \leq n \sum
\lambda_k^2$, which gives $d_{\text{eff}} \geq (\sum \lambda_k)^2 / (n \cdot
\max_k \lambda_k \cdot \sum_l \lambda_l)$ --- actually, the lower bound
follows from $\sum \lambda_k^2 \leq (\sum \lambda_k)^2$ being false in
general. The correct argument: $\sum \lambda_k^2 \leq (\max \lambda_k)(\sum
\lambda_k) \leq (\sum \lambda_k)^2$, so $d_{\text{eff}} \geq 1$.
\end{proof}

\begin{interpretation}[Is Culture Low-Dimensional?]
\label{interp:dimensionality}
The effective dimensionality $d_{\text{eff}}$ answers a fundamental empirical
question: \emph{how many independent axes of variation exist in a given
culture?}

Empirical evidence from political science, computational linguistics, and
social psychology consistently finds low effective dimensionality in cultural
data. Political attitudes are well-described by 1--3 dimensions (economic
left--right, social liberal--conservative, sometimes a third ``populist''
dimension). Semantic spaces extracted from text corpora are well-approximated
by 50--300 dimensions (orders of magnitude less than the vocabulary size).
Moral psychology identifies 5--6 ``moral foundations'' that span most of the
variation in moral judgment.

IDT~v2 provides a theoretical framework for understanding this low
dimensionality. If ideas combine additively and resonate via inner products,
then the resonance matrix is a Gram matrix, and its rank equals the dimension
of the meaning space actually occupied by the population's ideas. The
empirical finding of low dimensionality means that cultural populations
cluster in a low-dimensional subspace of the full Hilbert space ---
they use only a small fraction of the available ``meaning directions.''

This connects to the linguistic hypothesis of \emph{conceptual economy}: human
thought uses a small number of fundamental dimensions, and complex ideas are
compositions of these basic dimensions. The eigenvalues of the resonance matrix
quantify this economy: the rate of eigenvalue decay reveals how quickly
the ``explanatory power'' of each successive dimension diminishes.
\end{interpretation}


\section{Dynamics: How the Spectrum Evolves}

\begin{definition}[Spectral Evolution]
\label{def:spectral-evolution}
In a \textbf{dynamic cultural population}, each agent updates their idea over
time. At time $t$, the resonance matrix is $\RMat(t)$ with eigenvalues
$\lambda_1(t) \geq \lambda_2(t) \geq \cdots$.
\end{definition}

\begin{proposition}[Agreement Increases Under Averaging]
\label{prop:averaging}
If at each time step, each agent updates their idea to the average of their
neighbors' ideas (DeGroot learning):
\[
  \embed{a_i(t+1)} = \frac{1}{|N(i)|} \sum_{j \in N(i)} \embed{a_j(t)},
\]
where $N(i)$ is the set of agent $i$'s neighbors, then the spectral gap
$\gamma(t) = \lambda_1(t) - \lambda_2(t)$ is non-decreasing. The culture
becomes ``more one-dimensional'' over time.
\end{proposition}

\begin{proof}[Proof sketch]
Under DeGroot averaging, the embedding matrix evolves as $\Phi(t+1) = \Phi(t)
\cdot W$, where $W$ is the row-stochastic weight matrix. The resonance matrix
becomes $\RMat(t+1) = W^\top \RMat(t) W$. The spectral properties of this
iterated matrix product are well-studied: if $W$ is connected and aperiodic,
the eigenvalues of $\RMat(t)$ converge to a matrix with rank 1 (all agents
converge to the same idea), and the spectral gap increases monotonically.
\end{proof}

\begin{interpretation}[Convergence to Consensus]
\label{interp:convergence}
Proposition~\ref{prop:averaging} formalizes the common observation that social
averaging leads to consensus: when agents update their beliefs by averaging
with their neighbors, the population's effective dimensionality decreases over
time. Eventually, all agents hold approximately the same idea (the ``average''
of the initial ideas, weighted by network centrality).

This is the spectral-theoretic version of the convergence results in
Chapter~\ref{ch:persuasion}: repeated mutual influence (iterated persuasion)
drives the population toward a one-dimensional consensus. The spectral gap
quantifies how quickly this convergence occurs: a large gap means fast
convergence; a small gap means slow convergence.

The result also reveals the \emph{cost} of consensus: dimensional reduction.
As the culture converges on a single direction, it loses the richness of its
initial multi-dimensionality. The diverse voices are averaged into a bland
consensus. This is the mathematical foundation of the ``tyranny of the
majority'' or the ``spiral of silence'': social averaging suppresses minority
views (eigenvalues $\lambda_2, \lambda_3, \ldots$) in favor of the dominant
mode ($\lambda_1$).

In a richer model, one might introduce \emph{innovation} (random perturbations
that increase dimensionality) to counteract the averaging process. The
interplay between averaging (which decreases dimensionality) and innovation
(which increases it) would determine the equilibrium dimensionality of the
culture --- a theme we develop in Part~III.
\end{interpretation}

\begin{example}[Spectral Evolution of a Polarizing Society]
\label{ex:polarizing}
Consider a society that starts with five agents in $\mathbb{R}^2$:
\begin{align*}
  \embed{a_1(0)} &= (2, 1), \quad \embed{a_2(0)} = (1, 2), \quad
  \embed{a_3(0)} = (0, 1), \\
  \embed{a_4(0)} &= (-1, 1), \quad \embed{a_5(0)} = (-2, 0).
\end{align*}
Initially, the eigenvalues of $\RMat$ are approximately $\lambda_1 \approx 11.5$
and $\lambda_2 \approx 4.5$, giving $d_{\text{eff}} \approx 1.6$ (nearly
two-dimensional) and $\gamma \approx 7.0$.

If agents 1--3 average with each other and agents 4--5 average with each other
(but the two groups do not interact), the society splits into two clusters:
agents 1--3 converge to approximately $(1, 1.3)$ and agents 4--5 converge to
approximately $(-1.5, 0.5)$. The resonance matrix becomes approximately
block-diagonal, with $\lambda_1 \approx 9.5$, $\lambda_2 \approx 5.0$. The
spectral gap \emph{decreases} from 7.0 to 4.5, and the effective
dimensionality \emph{increases} from 1.6 to 1.9. Polarization increases the
effective dimensionality and decreases the spectral gap --- the opposite of
what happens under global averaging.

This illustrates a key insight: \emph{local averaging within clusters increases
within-cluster agreement but increases between-cluster divergence.} The
culture becomes ``more two-dimensional'' as the two clusters drift apart.
Polarization is, in spectral terms, the growth of $\lambda_2$ relative to
$\lambda_1$.
\end{example}

\bigskip

This chapter concludes Part~II of the book. We have developed the spectral
theory of cultural agreement: the resonance matrix is a Gram matrix (always
positive semidefinite), its eigenvalues reveal the principal modes of
agreement, the spectral gap measures cultural cohesion, and spectral
clustering identifies cultural tribes. The eigenvector of the largest eigenvalue
is Durkheim's collective conscience; the effective dimensionality measures the
complexity of the cultural landscape; and the dynamics of spectral evolution
reveal the tension between consensus (dimensional reduction) and diversity
(dimensional expansion).

In Part~III, we will study the dynamics of meaning change in Hilbert space:
how ideas move, how populations evolve, how learning and innovation interact
with the geometric structure of meaning.
